\clearpage
\fancyhead[L]{\small DE LATITUDINE SPECULAE MANHEMIENSIS.}
\hypertarget{native-602-lower}{}
\begin{center}\small Göttingische gelehrte Anzeigen. Stück 76. Seiten 753..756. 12. Mai 1817.\\[3pt]\rule{.72\linewidth}{.3pt}\end{center}

\emph{De latitudine speculae Manhemiensis autore} H.\,C.\ \textsc{Schumacher},
\emph{astronomiae professore et regiae societatis scientiarum Havniensis
socio.} Kopenhagen 1816. Gedruckt bei J. F. Schulz. 56 Seiten im Großquartformat.

\begingroup\brokenpenalty=10000\relax
Die Sternwarte in Mannheim besitzt zur Messung der Zenitdistanzen einen Mauer\SourceLeaf{603}quadranten von \textsc{Bird} mit acht Fuß Radius, einen Neunfuß-Zenitsektor von \textsc{Sisson} und seit 1811 auch einen Dreifuß-Repetitionskreis von \textsc{Reichenbach}. Für die Aufstellung des letzten Instruments wurden jedoch bisher keine Vorkehrungen getroffen. Der Verfasser, von 1813 bis 1815 Direktor dieser Sternwarte, war daher auf die beiden ersten Instrumente beschränkt, um die Polhöhe der Sternwarte so genau zu bestimmen, wie es die vorhandenen Hilfsmittel erlaubten. Der Zenitsektor diente zur Bestimmung des Kollimationsfehlers des Mauerquadranten. Da der Mauerquadrant nach Süden ausgerichtet ist und nicht umgehängt werden kann, war eine eigenständige Bestimmung der Polhöhe anhand von Fixsternen unmöglich. Es blieb nur, die Deklinationen der beobachteten Fixsterne nach den zuverlässigsten und neuesten Bestimmungen anderer Astronomen zugrunde zu legen. Mit den 152 Beobachtungen, die vom 27. Januar 1814 bis zum 4. März 1815 am Zenitsektor angestellt wurden, bestimmte der Verfasser die mittleren, auf den Anfang des Jahres 1815 reduzierten Zenitdistanzen von achtzehn Sternen. Diese bedurften jedoch einer Korrektur, da die Prüfung ergab, dass die Gradteilungen auf dem Zenitsektor zu klein und die gemessenen Zenitdistanzen folglich zu groß waren.
Um das Gesetz dieser Korrekturen zu bestimmen, maß der Verfasser mit einem Stangenzirkel die Sehnen mehrerer Bögen auf dem Limbus des Sektors und zugleich die Entfernung vom Mittelpunkt der Bewegung.
Die Formel für die Korrektur, die Herr \textsc{Bessel} in Königsberg aus diesen Daten ableitete, lautet
\[
1''3892\,z
\]
wobei $z$ die Zenitdistanz in Grad bezeichnet. Sie gibt jene Messungen recht gut wieder, und die kleinen Differenzen lassen sich mit gutem Grund als Messfehler ansehen.
Der Verfasser hat indessen die Formel
\[
0''4225\left(z+\tfrac{1}{10}zz+\tfrac{1}{20}z^{3}\right)
\]
vorgezogen, die eine noch bessere Übereinstimmung zeigt.
Hierbei wird also angenommen, dass die Gradteilungen des Zenitsektors zu klein und außerdem unterschiedlich groß sind.
Genau genommen ist das zweite Glied hier jedoch unzulässig, weil der Nullpunkt alle Bögen halbiert, deren Sehnen gemessen wurden.
Auch hätte die Formel
\[
0''585\left(z+\tfrac{1}{28}z^{3}\right)
\]
eine ebenso gute Übereinstimmung ergeben.
Im Übrigen ist es für das Endergebnis fast völlig gleichgültig, welche dieser Korrekturformeln man verwendet. Die so gefundenen Zenitdistanzen dienten nun zur Bestimmung des Kollimationsfehlers des Mauerquadranten, mit dem dieselben Sterne beobachtet worden waren. Mit diesem Fehler wurden die Zenitdistanzen von 36 weiteren, am Mauerquadranten beobachteten Sternen berechnet. Der Vergleich mit den von \textsc{Piazzi}, \textsc{Pond} und \textsc{Oriani} bestimmten Deklinationen ergab für die Polhöhe der Mannheimer Sternwarte als Ergebnis aus 284 Beobachtungen $49^{\circ}\,29'\,13''5$. Von den 18 am Zenitsektor beobachteten Sternen selbst konnte der Verfasser nur neun auf dieselbe Weise verwenden, da ihm für die übrigen ebenso zuverlässige Bestimmungen der Deklination fehlten. Diese neun ergaben für ihn aus 87 Beobachtungen die Polhöhe $49^{\circ}\,29'\,14''2$. Der Rezensent hat die Beobachtungen der übrigen neun anhand der Deklinationsbestimmungen in \textsc{Piazzi}s neuem Katalog berechnet, den der Verfasser noch nicht benutzen konnte, und folgende Ergebnisse gefunden:
\begin{center}
\fontsize{9}{11}\selectfont
\setlength{\tabcolsep}{3pt}
\begin{tabular}{l|c|c||l|c|c}
 & Polhöhe & \begin{tabular}{c}Anzahl der\\Beobachtungen\end{tabular} & & Polhöhe & \begin{tabular}{c}Anzahl der\\Beobachtungen\end{tabular} \\
\hline
$\sigma$ Perseus & $49^{\circ}\,29'\,10''4$ & 8 & 725 Perseus & $49^{\circ}\,29'\,11''5$ & 9 \\
$\psi$ Perseus & 8.3 & 9 & 7 Perseus & 8.2 & 8 \\
$\chi$ Gr. Bär & 8 3 & 5 & $F$ Gr. Bär & 9.8 & 5 \\
$\theta$ Perseus & 15.2 & 5 & $\gamma$ Perseus & 8.2 & 13 \\
$\iota$ Gr. Bär & 8.0 & 5 & \multicolumn{3}{|l}{} \\
\end{tabular}
\end{center}
\noindent\SourceLeaf{604}Für den vom Verfasser als 725 Perseus bezeichneten Stern, dessen Bezeichnung sich nicht finden ließ, wurde der Eintrag II, 253 in \textsc{Piazzi}s Katalog angenommen. Diese 65 Beobachtungen ergeben also im Mittel $49^\circ\,29'\,9''8$. Folglich ergeben alle 152 Beobachtungen am Zenitsektor im Mittel $49^\circ\,29'\,12''3$ und damit sämtliche 436 Fixsternbeobachtungen im Mittel $49^\circ\,29'\,13''1$. Die 66 Sonnenbeobachtungen, die der Verfasser im Jahr 1814 anstellte, ergeben im Mittel nur einige Zehntel einer Sekunde mehr, können aber aus einigen örtlichen Gründen nicht dieselbe Zuverlässigkeit beanspruchen. Dieses Ergebnis aus den Beobachtungen des Verfassers wird nun als das sicherste gelten müssen, bis der Kreis von \textsc{Reichenbach} aufgestellt und zu der einer solchen Sternwarte eigentlich allein würdigen neuen Bestimmung der Polhöhe anhand von Zirkumpolarsternen verwendet worden ist\enspace Übrigens stimmt das obige Ergebnis zwar genau mit der letzten Angabe von \textsc{Barry} überein, doch hatte dieser Astronom nicht bekannt gemacht, wie er zu seiner Bestimmung gelangt war. Da sich deshalb der Grad ihrer Genauigkeit nicht beurteilen ließ, verdient Herr \textsc{Schumacher} für die vollständige Veröffentlichung seiner Beobachtungen den Dank aller Astronomen.
\par\clearpage\endgroup
